\begin{table}[htbp]
\centering
\footnotesize
\setlength{\tabcolsep}{3pt}
\caption{流感规则窗口（$S=4$）的权重比 $q_1$、$q_2$ 与纯报告噪声零模型的比较。}
\label{tab:s6}
\begin{tabular}{rrrcccc}
\toprule
$\mathrm{CV}^2$ & $R$ & 有效窗口 & $q_1$ 均值（MCSE） & $q_1$ 95\% 范围 & $q_2$ 均值（MCSE） & $q_2$ 95\% 范围 \\
\midrule
0.0036 & 300 & 262 & 0.760 (0.001) & [0.729, 0.786] & 1.464 (0.002) & [1.417, 1.525] \\
0.0143 & 300 & 231 & 0.900 (0.002) & [0.853, 0.955] & 1.282 (0.003) & [1.211, 1.356] \\
0.0287 & 300 & 238 & 0.990 (0.002) & [0.913, 1.065] & 1.167 (0.003) & [1.080, 1.267] \\
0.0574 & 300 & 208 & 1.067 (0.003) & [0.981, 1.158] & 1.055 (0.003) & [0.959, 1.161] \\
\midrule
\multicolumn{3}{l}{零模型范围（各格均值的极差）} & \multicolumn{2}{c}{0.760--1.067} & \multicolumn{2}{c}{1.055--1.464} \\
\multicolumn{3}{l}{敏感性：各季窗口均 $\ge4$ 周} & \multicolumn{2}{c}{0.762--1.077} & \multicolumn{2}{c}{1.051--1.462} \\
\multicolumn{3}{l}{观测（流感，$N=1756$）} & \multicolumn{2}{c}{0.660 [0.562, 0.792]} & \multicolumn{2}{c}{1.597 [1.384, 1.845]} \\
\multicolumn{3}{l}{观测落在零模型范围外的 $p$} & \multicolumn{2}{c}{0.120} & \multicolumn{2}{c}{0.217} \\
\bottomrule
\end{tabular}
\par\smallskip
\begin{minipage}{0.95\linewidth}\scriptsize $q_1=a_{\text{单周}}/a_{\text{原权重}}$，$q_2=a_{\text{等权}}/a_{\text{原权重}}$；各权重各自取 $D\ge5$ 且四周覆盖率为正的行。零模型条件于平滑的全国季节曲线与各州份额，只加入对数正态报告率（方差 $\mathrm{CV}^2$）与泊松计数噪声。每条模拟路径都按观测数据的规则重新选择窗口。任一季无法形成窗口的路径记为无效，不参与汇总；无效路径主要来自 2025-26 季起点处增长已低于 1 的情形。95\% 范围为模拟分布的 2.5\% 与 97.5\% 分位数，其分位点 CDF 尺度的 Monte Carlo 误差约 $\sqrt{0.025\cdot0.975/R_{\text{有效}}}\approx0.01$。该零模型不含传播增长、代际间隔与入院延迟。$p$ 为观测自助分布（$B=2000$，与 $G$ 规则窗口同一重抽流）落在零模型范围一侧比例的两倍，未做多重校正。\end{minipage}
\end{table}
